% paper50-semantic-centers.tex — Semantic Center Detection: Finding the Most
% Important Code to Verify First
% Paper 50 of the Judgment Geometry series.
% Compile: pdflatex -interaction=nonstopmode paper50-semantic-centers.tex
\documentclass[11pt]{article}
\usepackage{lmodern}
% jugeo-common.tex — shared preamble for all Judgment Geometry papers
% Include via: % jugeo-common.tex — shared preamble for all Judgment Geometry papers
% Include via: % jugeo-common.tex — shared preamble for all Judgment Geometry papers
% Include via: \input{jugeo-common}

% ── Packages ────────────────────────────────────────────────────────────
\usepackage{amsmath,amssymb,amsthm,mathtools}
\usepackage{tikz-cd}
\usepackage{listings}
\usepackage{xcolor}
\usepackage{xspace}
\usepackage{booktabs}
\usepackage{multirow}
\usepackage{enumitem}
\usepackage[expansion=false]{microtype}
\usepackage{hyperref}
\usepackage{cleveref}
% \usepackage{thmtools}  % disabled: not available in this TeX installation
\usepackage{float}
\usepackage{caption}
\usepackage{subcaption}
\usepackage{tabularx}
\usepackage{stmaryrd}       % \llbracket, \rrbracket
\usepackage{bbm}            % \mathbbm{1}

% ── Page geometry ───────────────────────────────────────────────────────
\usepackage[margin=1in]{geometry}

% ── Theorem environments ────────────────────────────────────────────────
\theoremstyle{plain}
\newtheorem{theorem}{Theorem}[section]
\newtheorem{lemma}[theorem]{Lemma}
\newtheorem{proposition}[theorem]{Proposition}
\newtheorem{corollary}[theorem]{Corollary}

\theoremstyle{definition}
\newtheorem{definition}[theorem]{Definition}
\newtheorem{example}[theorem]{Example}
\newtheorem{construction}[theorem]{Construction}

\theoremstyle{remark}
\newtheorem{remark}[theorem]{Remark}
\newtheorem{notation}[theorem]{Notation}

% ── Python listings style ──────────────────────────────────────────────
\definecolor{jugeoBlue}{HTML}{2563EB}
\definecolor{jugeoGreen}{HTML}{059669}
\definecolor{jugeoRed}{HTML}{DC2626}
\definecolor{jugeoPurple}{HTML}{7C3AED}
\definecolor{jugeoGray}{HTML}{6B7280}
\definecolor{jugeoBg}{HTML}{F8FAFC}

\lstdefinestyle{jugeo-python}{
  language=Python,
  basicstyle=\ttfamily\small,
  keywordstyle=\color{jugeoBlue}\bfseries,
  stringstyle=\color{jugeoGreen},
  commentstyle=\color{jugeoGray}\itshape,
  numberstyle=\tiny\color{jugeoGray},
  backgroundcolor=\color{jugeoBg},
  numbers=left,
  numbersep=6pt,
  frame=single,
  framerule=0.4pt,
  rulecolor=\color{jugeoGray!40},
  breaklines=true,
  breakatwhitespace=true,
  showstringspaces=false,
  tabsize=4,
  xleftmargin=1.5em,
  framexleftmargin=1.5em,
  morekeywords={dataclass,frozen,field,Enum,IntEnum,auto,Any,
                Mapping,Sequence,Optional,match,case},
  literate={→}{{$\to$}}1 {←}{{$\leftarrow$}}1
           {≤}{{$\leq$}}1 {≥}{{$\geq$}}1 {≠}{{$\neq$}}1
           {∈}{{$\in$}}1 {∀}{{$\forall$}}1 {∃}{{$\exists$}}1
           {⊕}{{$\oplus$}}1 {⊖}{{$\ominus$}}1,
  escapeinside={(*@}{@*)},
}
\lstset{style=jugeo-python}

\lstdefinestyle{jugeo-output}{
  basicstyle=\ttfamily\small,
  backgroundcolor=\color{jugeoBg},
  frame=single,
  framerule=0.4pt,
  rulecolor=\color{jugeoGray!40},
  breaklines=true,
  showstringspaces=false,
  xleftmargin=1.5em,
  framexleftmargin=0.5em,
  numbers=none,
}

% ── JuGeo-specific macros ──────────────────────────────────────────────

% Judgment 8-tuple
\newcommand{\J}{\mathcal{J}}
\newcommand{\Jud}[8]{(#1,\,#2,\,#3,\,#4,\,#5,\,#6,\,#7,\,#8)}
\newcommand{\JudShort}[1]{\J_{#1}}

% Components
\newcommand{\coord}{\mathsf{c}}            % coordinate
\newcommand{\prop}{\varphi}                 % proposition
\newcommand{\carrier}{\mathcal{A}}          % carrier type
\newcommand{\evid}{\mathcal{E}}             % evidence bundle
\newcommand{\oblig}{\mathcal{O}}            % obligations
\newcommand{\obstr}{\mathcal{B}}            % obstructions
\newcommand{\trust}{\mathcal{T}}            % trust annotation
\newcommand{\prov}{\Pi}                     % provenance

% Trust algebra
\newcommand{\Talg}{\mathfrak{T}}            % trust ordered algebra
\newcommand{\Eadm}{\mathcal{E}_{\mathrm{adm}}}
\newcommand{\tleq}{\preceq}                 % trust ordering
\newcommand{\tjoin}{\oplus}                 % conservative join
\newcommand{\tatten}{\ominus}               % attenuation
\newcommand{\tpromote}{\uparrow_\pi}        % promotion
\newcommand{\tdemote}{\downarrow_\chi}       % demotion

% Trust tiers
\newcommand{\Tcontra}{\ensuremath{\mathsf{CONTRADICTED}}}
\newcommand{\Tunver}{\ensuremath{\mathsf{UNVERIFIED}}}
\newcommand{\Tcopilot}{\ensuremath{\mathsf{COPILOT}}}
\newcommand{\Toracle}{\ensuremath{\mathsf{ORACLE}}}
\newcommand{\Truntime}{\ensuremath{\mathsf{RUNTIME}}}
\newcommand{\Tsolver}{\ensuremath{\mathsf{SOLVER}}}
\newcommand{\Tproof}{\ensuremath{\mathsf{PROOF}}}

% Site and geometry
\newcommand{\Site}{\mathcal{S}}             % semantic site
\newcommand{\Cov}{\mathrm{Cov}}            % covering family
\newcommand{\Ob}{\mathrm{Ob}}              % objects
\newcommand{\Mor}{\mathrm{Mor}}            % morphisms
\newcommand{\res}{\mathrm{res}}            % restriction
\newcommand{\Cech}{\check{C}}              % Čech complex
\newcommand{\Hcoh}[1]{H^{#1}}             % cohomology group

% Descent
\newcommand{\descent}{\mathrm{desc}}
\newcommand{\glue}{\mathrm{glue}}
\newcommand{\overlap}[2]{#1 \cap #2}

% Semantic moves
\newcommand{\VERIFY}{\textsc{Verify}}
\newcommand{\CONSTRUCT}{\textsc{Construct}}
\newcommand{\NORMALIZE}{\textsc{Normalize}}
\newcommand{\GLUE}{\textsc{Glue}}
\newcommand{\RESTRICT}{\textsc{Restrict}}
\newcommand{\TRANSPORT}{\textsc{Transport}}
\newcommand{\REPAIR}{\textsc{Repair}}
\newcommand{\PROMOTE}{\textsc{Promote}}
\newcommand{\CHALLENGE}{\textsc{Challenge}}
\newcommand{\ESCALATE}{\textsc{Escalate}}

% Fragments
\newcommand{\QFLIA}{\texttt{QF\_LIA}}
\newcommand{\QFLRA}{\texttt{QF\_LRA}}
\newcommand{\QFBV}{\texttt{QF\_BV}}
\newcommand{\QFUF}{\texttt{QF\_UF}}

% Misc
\newcommand{\Python}{\textsc{Python}}
\newcommand{\JuGeo}{\textsc{JuGeo}}
\newcommand{\LEAN}{\textsc{Lean}\,4}
\newcommand{\Fstar}{F$^{\star}$}
\newcommand{\Coq}{\textsc{Coq}}
\newcommand{\Dafny}{\textsc{Dafny}}

% Common shortcuts
\newcommand{\ie}{i.e.\@\xspace}
\newcommand{\eg}{e.g.\@\xspace}
\newcommand{\cf}{cf.\@\xspace}
\newcommand{\wrt}{w.r.t.\@\xspace}

% Evaluation shortcuts
\newcommand{\numcases}{300}
\newcommand{\numspec}{100}
\newcommand{\numeq}{100}
\newcommand{\numbug}{100}
\newcommand{\medtime}{6.5\,ms}
\newcommand{\totaltime}{1.949\,s}


% ── Packages ────────────────────────────────────────────────────────────
\usepackage{amsmath,amssymb,amsthm,mathtools}
\usepackage{tikz-cd}
\usepackage{listings}
\usepackage{xcolor}
\usepackage{xspace}
\usepackage{booktabs}
\usepackage{multirow}
\usepackage{enumitem}
\usepackage[expansion=false]{microtype}
\usepackage{hyperref}
\usepackage{cleveref}
% \usepackage{thmtools}  % disabled: not available in this TeX installation
\usepackage{float}
\usepackage{caption}
\usepackage{subcaption}
\usepackage{tabularx}
\usepackage{stmaryrd}       % \llbracket, \rrbracket
\usepackage{bbm}            % \mathbbm{1}

% ── Page geometry ───────────────────────────────────────────────────────
\usepackage[margin=1in]{geometry}

% ── Theorem environments ────────────────────────────────────────────────
\theoremstyle{plain}
\newtheorem{theorem}{Theorem}[section]
\newtheorem{lemma}[theorem]{Lemma}
\newtheorem{proposition}[theorem]{Proposition}
\newtheorem{corollary}[theorem]{Corollary}

\theoremstyle{definition}
\newtheorem{definition}[theorem]{Definition}
\newtheorem{example}[theorem]{Example}
\newtheorem{construction}[theorem]{Construction}

\theoremstyle{remark}
\newtheorem{remark}[theorem]{Remark}
\newtheorem{notation}[theorem]{Notation}

% ── Python listings style ──────────────────────────────────────────────
\definecolor{jugeoBlue}{HTML}{2563EB}
\definecolor{jugeoGreen}{HTML}{059669}
\definecolor{jugeoRed}{HTML}{DC2626}
\definecolor{jugeoPurple}{HTML}{7C3AED}
\definecolor{jugeoGray}{HTML}{6B7280}
\definecolor{jugeoBg}{HTML}{F8FAFC}

\lstdefinestyle{jugeo-python}{
  language=Python,
  basicstyle=\ttfamily\small,
  keywordstyle=\color{jugeoBlue}\bfseries,
  stringstyle=\color{jugeoGreen},
  commentstyle=\color{jugeoGray}\itshape,
  numberstyle=\tiny\color{jugeoGray},
  backgroundcolor=\color{jugeoBg},
  numbers=left,
  numbersep=6pt,
  frame=single,
  framerule=0.4pt,
  rulecolor=\color{jugeoGray!40},
  breaklines=true,
  breakatwhitespace=true,
  showstringspaces=false,
  tabsize=4,
  xleftmargin=1.5em,
  framexleftmargin=1.5em,
  morekeywords={dataclass,frozen,field,Enum,IntEnum,auto,Any,
                Mapping,Sequence,Optional,match,case},
  literate={→}{{$\to$}}1 {←}{{$\leftarrow$}}1
           {≤}{{$\leq$}}1 {≥}{{$\geq$}}1 {≠}{{$\neq$}}1
           {∈}{{$\in$}}1 {∀}{{$\forall$}}1 {∃}{{$\exists$}}1
           {⊕}{{$\oplus$}}1 {⊖}{{$\ominus$}}1,
  escapeinside={(*@}{@*)},
}
\lstset{style=jugeo-python}

\lstdefinestyle{jugeo-output}{
  basicstyle=\ttfamily\small,
  backgroundcolor=\color{jugeoBg},
  frame=single,
  framerule=0.4pt,
  rulecolor=\color{jugeoGray!40},
  breaklines=true,
  showstringspaces=false,
  xleftmargin=1.5em,
  framexleftmargin=0.5em,
  numbers=none,
}

% ── JuGeo-specific macros ──────────────────────────────────────────────

% Judgment 8-tuple
\newcommand{\J}{\mathcal{J}}
\newcommand{\Jud}[8]{(#1,\,#2,\,#3,\,#4,\,#5,\,#6,\,#7,\,#8)}
\newcommand{\JudShort}[1]{\J_{#1}}

% Components
\newcommand{\coord}{\mathsf{c}}            % coordinate
\newcommand{\prop}{\varphi}                 % proposition
\newcommand{\carrier}{\mathcal{A}}          % carrier type
\newcommand{\evid}{\mathcal{E}}             % evidence bundle
\newcommand{\oblig}{\mathcal{O}}            % obligations
\newcommand{\obstr}{\mathcal{B}}            % obstructions
\newcommand{\trust}{\mathcal{T}}            % trust annotation
\newcommand{\prov}{\Pi}                     % provenance

% Trust algebra
\newcommand{\Talg}{\mathfrak{T}}            % trust ordered algebra
\newcommand{\Eadm}{\mathcal{E}_{\mathrm{adm}}}
\newcommand{\tleq}{\preceq}                 % trust ordering
\newcommand{\tjoin}{\oplus}                 % conservative join
\newcommand{\tatten}{\ominus}               % attenuation
\newcommand{\tpromote}{\uparrow_\pi}        % promotion
\newcommand{\tdemote}{\downarrow_\chi}       % demotion

% Trust tiers
\newcommand{\Tcontra}{\ensuremath{\mathsf{CONTRADICTED}}}
\newcommand{\Tunver}{\ensuremath{\mathsf{UNVERIFIED}}}
\newcommand{\Tcopilot}{\ensuremath{\mathsf{COPILOT}}}
\newcommand{\Toracle}{\ensuremath{\mathsf{ORACLE}}}
\newcommand{\Truntime}{\ensuremath{\mathsf{RUNTIME}}}
\newcommand{\Tsolver}{\ensuremath{\mathsf{SOLVER}}}
\newcommand{\Tproof}{\ensuremath{\mathsf{PROOF}}}

% Site and geometry
\newcommand{\Site}{\mathcal{S}}             % semantic site
\newcommand{\Cov}{\mathrm{Cov}}            % covering family
\newcommand{\Ob}{\mathrm{Ob}}              % objects
\newcommand{\Mor}{\mathrm{Mor}}            % morphisms
\newcommand{\res}{\mathrm{res}}            % restriction
\newcommand{\Cech}{\check{C}}              % Čech complex
\newcommand{\Hcoh}[1]{H^{#1}}             % cohomology group

% Descent
\newcommand{\descent}{\mathrm{desc}}
\newcommand{\glue}{\mathrm{glue}}
\newcommand{\overlap}[2]{#1 \cap #2}

% Semantic moves
\newcommand{\VERIFY}{\textsc{Verify}}
\newcommand{\CONSTRUCT}{\textsc{Construct}}
\newcommand{\NORMALIZE}{\textsc{Normalize}}
\newcommand{\GLUE}{\textsc{Glue}}
\newcommand{\RESTRICT}{\textsc{Restrict}}
\newcommand{\TRANSPORT}{\textsc{Transport}}
\newcommand{\REPAIR}{\textsc{Repair}}
\newcommand{\PROMOTE}{\textsc{Promote}}
\newcommand{\CHALLENGE}{\textsc{Challenge}}
\newcommand{\ESCALATE}{\textsc{Escalate}}

% Fragments
\newcommand{\QFLIA}{\texttt{QF\_LIA}}
\newcommand{\QFLRA}{\texttt{QF\_LRA}}
\newcommand{\QFBV}{\texttt{QF\_BV}}
\newcommand{\QFUF}{\texttt{QF\_UF}}

% Misc
\newcommand{\Python}{\textsc{Python}}
\newcommand{\JuGeo}{\textsc{JuGeo}}
\newcommand{\LEAN}{\textsc{Lean}\,4}
\newcommand{\Fstar}{F$^{\star}$}
\newcommand{\Coq}{\textsc{Coq}}
\newcommand{\Dafny}{\textsc{Dafny}}

% Common shortcuts
\newcommand{\ie}{i.e.\@\xspace}
\newcommand{\eg}{e.g.\@\xspace}
\newcommand{\cf}{cf.\@\xspace}
\newcommand{\wrt}{w.r.t.\@\xspace}

% Evaluation shortcuts
\newcommand{\numcases}{300}
\newcommand{\numspec}{100}
\newcommand{\numeq}{100}
\newcommand{\numbug}{100}
\newcommand{\medtime}{6.5\,ms}
\newcommand{\totaltime}{1.949\,s}


% ── Packages ────────────────────────────────────────────────────────────
\usepackage{amsmath,amssymb,amsthm,mathtools}
\usepackage{tikz-cd}
\usepackage{listings}
\usepackage{xcolor}
\usepackage{xspace}
\usepackage{booktabs}
\usepackage{multirow}
\usepackage{enumitem}
\usepackage[expansion=false]{microtype}
\usepackage{hyperref}
\usepackage{cleveref}
% \usepackage{thmtools}  % disabled: not available in this TeX installation
\usepackage{float}
\usepackage{caption}
\usepackage{subcaption}
\usepackage{tabularx}
\usepackage{stmaryrd}       % \llbracket, \rrbracket
\usepackage{bbm}            % \mathbbm{1}

% ── Page geometry ───────────────────────────────────────────────────────
\usepackage[margin=1in]{geometry}

% ── Theorem environments ────────────────────────────────────────────────
\theoremstyle{plain}
\newtheorem{theorem}{Theorem}[section]
\newtheorem{lemma}[theorem]{Lemma}
\newtheorem{proposition}[theorem]{Proposition}
\newtheorem{corollary}[theorem]{Corollary}

\theoremstyle{definition}
\newtheorem{definition}[theorem]{Definition}
\newtheorem{example}[theorem]{Example}
\newtheorem{construction}[theorem]{Construction}

\theoremstyle{remark}
\newtheorem{remark}[theorem]{Remark}
\newtheorem{notation}[theorem]{Notation}

% ── Python listings style ──────────────────────────────────────────────
\definecolor{jugeoBlue}{HTML}{2563EB}
\definecolor{jugeoGreen}{HTML}{059669}
\definecolor{jugeoRed}{HTML}{DC2626}
\definecolor{jugeoPurple}{HTML}{7C3AED}
\definecolor{jugeoGray}{HTML}{6B7280}
\definecolor{jugeoBg}{HTML}{F8FAFC}

\lstdefinestyle{jugeo-python}{
  language=Python,
  basicstyle=\ttfamily\small,
  keywordstyle=\color{jugeoBlue}\bfseries,
  stringstyle=\color{jugeoGreen},
  commentstyle=\color{jugeoGray}\itshape,
  numberstyle=\tiny\color{jugeoGray},
  backgroundcolor=\color{jugeoBg},
  numbers=left,
  numbersep=6pt,
  frame=single,
  framerule=0.4pt,
  rulecolor=\color{jugeoGray!40},
  breaklines=true,
  breakatwhitespace=true,
  showstringspaces=false,
  tabsize=4,
  xleftmargin=1.5em,
  framexleftmargin=1.5em,
  morekeywords={dataclass,frozen,field,Enum,IntEnum,auto,Any,
                Mapping,Sequence,Optional,match,case},
  literate={→}{{$\to$}}1 {←}{{$\leftarrow$}}1
           {≤}{{$\leq$}}1 {≥}{{$\geq$}}1 {≠}{{$\neq$}}1
           {∈}{{$\in$}}1 {∀}{{$\forall$}}1 {∃}{{$\exists$}}1
           {⊕}{{$\oplus$}}1 {⊖}{{$\ominus$}}1,
  escapeinside={(*@}{@*)},
}
\lstset{style=jugeo-python}

\lstdefinestyle{jugeo-output}{
  basicstyle=\ttfamily\small,
  backgroundcolor=\color{jugeoBg},
  frame=single,
  framerule=0.4pt,
  rulecolor=\color{jugeoGray!40},
  breaklines=true,
  showstringspaces=false,
  xleftmargin=1.5em,
  framexleftmargin=0.5em,
  numbers=none,
}

% ── JuGeo-specific macros ──────────────────────────────────────────────

% Judgment 8-tuple
\newcommand{\J}{\mathcal{J}}
\newcommand{\Jud}[8]{(#1,\,#2,\,#3,\,#4,\,#5,\,#6,\,#7,\,#8)}
\newcommand{\JudShort}[1]{\J_{#1}}

% Components
\newcommand{\coord}{\mathsf{c}}            % coordinate
\newcommand{\prop}{\varphi}                 % proposition
\newcommand{\carrier}{\mathcal{A}}          % carrier type
\newcommand{\evid}{\mathcal{E}}             % evidence bundle
\newcommand{\oblig}{\mathcal{O}}            % obligations
\newcommand{\obstr}{\mathcal{B}}            % obstructions
\newcommand{\trust}{\mathcal{T}}            % trust annotation
\newcommand{\prov}{\Pi}                     % provenance

% Trust algebra
\newcommand{\Talg}{\mathfrak{T}}            % trust ordered algebra
\newcommand{\Eadm}{\mathcal{E}_{\mathrm{adm}}}
\newcommand{\tleq}{\preceq}                 % trust ordering
\newcommand{\tjoin}{\oplus}                 % conservative join
\newcommand{\tatten}{\ominus}               % attenuation
\newcommand{\tpromote}{\uparrow_\pi}        % promotion
\newcommand{\tdemote}{\downarrow_\chi}       % demotion

% Trust tiers
\newcommand{\Tcontra}{\ensuremath{\mathsf{CONTRADICTED}}}
\newcommand{\Tunver}{\ensuremath{\mathsf{UNVERIFIED}}}
\newcommand{\Tcopilot}{\ensuremath{\mathsf{COPILOT}}}
\newcommand{\Toracle}{\ensuremath{\mathsf{ORACLE}}}
\newcommand{\Truntime}{\ensuremath{\mathsf{RUNTIME}}}
\newcommand{\Tsolver}{\ensuremath{\mathsf{SOLVER}}}
\newcommand{\Tproof}{\ensuremath{\mathsf{PROOF}}}

% Site and geometry
\newcommand{\Site}{\mathcal{S}}             % semantic site
\newcommand{\Cov}{\mathrm{Cov}}            % covering family
\newcommand{\Ob}{\mathrm{Ob}}              % objects
\newcommand{\Mor}{\mathrm{Mor}}            % morphisms
\newcommand{\res}{\mathrm{res}}            % restriction
\newcommand{\Cech}{\check{C}}              % Čech complex
\newcommand{\Hcoh}[1]{H^{#1}}             % cohomology group

% Descent
\newcommand{\descent}{\mathrm{desc}}
\newcommand{\glue}{\mathrm{glue}}
\newcommand{\overlap}[2]{#1 \cap #2}

% Semantic moves
\newcommand{\VERIFY}{\textsc{Verify}}
\newcommand{\CONSTRUCT}{\textsc{Construct}}
\newcommand{\NORMALIZE}{\textsc{Normalize}}
\newcommand{\GLUE}{\textsc{Glue}}
\newcommand{\RESTRICT}{\textsc{Restrict}}
\newcommand{\TRANSPORT}{\textsc{Transport}}
\newcommand{\REPAIR}{\textsc{Repair}}
\newcommand{\PROMOTE}{\textsc{Promote}}
\newcommand{\CHALLENGE}{\textsc{Challenge}}
\newcommand{\ESCALATE}{\textsc{Escalate}}

% Fragments
\newcommand{\QFLIA}{\texttt{QF\_LIA}}
\newcommand{\QFLRA}{\texttt{QF\_LRA}}
\newcommand{\QFBV}{\texttt{QF\_BV}}
\newcommand{\QFUF}{\texttt{QF\_UF}}

% Misc
\newcommand{\Python}{\textsc{Python}}
\newcommand{\JuGeo}{\textsc{JuGeo}}
\newcommand{\LEAN}{\textsc{Lean}\,4}
\newcommand{\Fstar}{F$^{\star}$}
\newcommand{\Coq}{\textsc{Coq}}
\newcommand{\Dafny}{\textsc{Dafny}}

% Common shortcuts
\newcommand{\ie}{i.e.\@\xspace}
\newcommand{\eg}{e.g.\@\xspace}
\newcommand{\cf}{cf.\@\xspace}
\newcommand{\wrt}{w.r.t.\@\xspace}

% Evaluation shortcuts
\newcommand{\numcases}{300}
\newcommand{\numspec}{100}
\newcommand{\numeq}{100}
\newcommand{\numbug}{100}
\newcommand{\medtime}{6.5\,ms}
\newcommand{\totaltime}{1.949\,s}

% experiment-data.tex — AUTO-GENERATED from real jugeo CLI runs
% DO NOT EDIT — regenerate with: python3 experiments/run_universal_metrics.py
% Generated from 30 programs across 6 domains

\newcommand{\expTotalPrograms}{30}
\newcommand{\expVerified}{30}
\newcommand{\expAccuracy}{100.0\%}
\newcommand{\expCoordsMin}{2}
\newcommand{\expCoordsMax}{10}
\newcommand{\expCoordsMean}{5.2}
\newcommand{\expCoordsMedian}{5.5}
\newcommand{\expPropsMin}{4}
\newcommand{\expPropsMax}{20}
\newcommand{\expPropsMean}{10.4}
\newcommand{\expPropsSum}{311}
\newcommand{\expPropsOkSum}{310}
\newcommand{\expTimeMin}{3.506\,s}
\newcommand{\expTimeMax}{6.714\,s}
\newcommand{\expTimeMean}{4.64\,s}
\newcommand{\expTimeMedian}{4.508\,s}
\newcommand{\expTimeTotal}{139.204\,s}
\newcommand{\expObstructionTotal}{0}
\newcommand{\expHOneZero}{30}
\newcommand{\expDomainCount}{6}
\newcommand{\expTrustCOPILOTSUGGESTED}{30}
\newcommand{\expDomainDsCount}{5}
\newcommand{\expDomainDsVerified}{5}
\newcommand{\expDomainDsAvgCoords}{7.6}
\newcommand{\expDomainDsAvgProps}{15.2}
\newcommand{\expDomainMathCount}{5}
\newcommand{\expDomainMathVerified}{5}
\newcommand{\expDomainMathAvgCoords}{4.8}
\newcommand{\expDomainMathAvgProps}{9.4}
\newcommand{\expDomainSmCount}{5}
\newcommand{\expDomainSmVerified}{5}
\newcommand{\expDomainSmAvgCoords}{7.6}
\newcommand{\expDomainSmAvgProps}{15.2}
\newcommand{\expDomainSortCount}{5}
\newcommand{\expDomainSortVerified}{5}
\newcommand{\expDomainSortAvgCoords}{2.4}
\newcommand{\expDomainSortAvgProps}{4.8}
\newcommand{\expDomainStrCount}{5}
\newcommand{\expDomainStrVerified}{5}
\newcommand{\expDomainStrAvgCoords}{3.2}
\newcommand{\expDomainStrAvgProps}{6.4}
\newcommand{\expDomainWebCount}{5}
\newcommand{\expDomainWebVerified}{5}
\newcommand{\expDomainWebAvgCoords}{5.6}
\newcommand{\expDomainWebAvgProps}{11.2}

% subsystem-data.tex — AUTO-GENERATED from real jugeo subsystem experiments
% DO NOT EDIT — regenerate with: python3 experiments/run_subsystem_experiments.py

\newcommand{\subCliTotal}{10}
\newcommand{\subCliVerified}{9}
\newcommand{\subCliAccuracy}{90.0\%}
\newcommand{\subCliMeanTime}{4.132\,s}
\newcommand{\subCliMinTime}{3.472\,s}
\newcommand{\subCliMaxTime}{5.372\,s}
\newcommand{\subCliTotalCoords}{54}
\newcommand{\subCliTotalProps}{107}
\newcommand{\subCliTotalPropsOk}{106}
\newcommand{\subSiteCalls}{90}
\newcommand{\subSiteSuccessful}{69}
\newcommand{\subSiteSuccessRate}{76.7\%}
\newcommand{\subSiteMeanTime}{0.0001\,s}
\newcommand{\subSiteTotalTime}{0.005\,s}
\newcommand{\subSpecificationsatisfactionSmallTime}{0.0003\,s}
\newcommand{\subSpecificationsatisfactionMediumTime}{0.0001\,s}
\newcommand{\subSpecificationsatisfactionLargeTime}{0.0001\,s}
\newcommand{\subInhabitantfleetSmallTime}{0.0\,s}
\newcommand{\subInhabitantfleetMediumTime}{0.0\,s}
\newcommand{\subInhabitantfleetLargeTime}{0.0\,s}
\newcommand{\subTheoremecologySmallTime}{0.0\,s}
\newcommand{\subTheoremecologyMediumTime}{0.0\,s}
\newcommand{\subTheoremecologyLargeTime}{0.0\,s}
\newcommand{\subStatespaceexplorationSmallTime}{0.0\,s}
\newcommand{\subStatespaceexplorationMediumTime}{0.0\,s}
\newcommand{\subStatespaceexplorationLargeTime}{0.0\,s}
\newcommand{\subRepairsemanticsSmallTime}{0.0\,s}
\newcommand{\subRepairsemanticsMediumTime}{0.0\,s}
\newcommand{\subRepairsemanticsLargeTime}{0.0\,s}
\newcommand{\subReplaygluingSmallTime}{0.0\,s}
\newcommand{\subReplaygluingMediumTime}{0.0\,s}
\newcommand{\subReplaygluingLargeTime}{0.0\,s}
\newcommand{\subSemanticfuturesSmallTime}{0.0\,s}
\newcommand{\subSemanticfuturesMediumTime}{0.0\,s}
\newcommand{\subSemanticfuturesLargeTime}{0.0\,s}
\newcommand{\subHypercovertreatySmallTime}{0.0\,s}
\newcommand{\subHypercovertreatyMediumTime}{0.0\,s}
\newcommand{\subHypercovertreatyLargeTime}{0.0\,s}
\newcommand{\subEncodeforsolverSmallTime}{0.0026\,s}
\newcommand{\subEncodeforsolverMediumTime}{0.0002\,s}
\newcommand{\subEncodeforsolverLargeTime}{0.0001\,s}
\newcommand{\subInterfaceroutingSmallTime}{0.0\,s}
\newcommand{\subInterfaceroutingMediumTime}{0.0\,s}
\newcommand{\subInterfaceroutingLargeTime}{0.0\,s}
\newcommand{\subOrchestrateverificationSmallTime}{0.0002\,s}
\newcommand{\subOrchestrateverificationMediumTime}{0.0\,s}
\newcommand{\subOrchestrateverificationLargeTime}{0.0\,s}
\newcommand{\subRunfulldescentSmallTime}{0.0\,s}
\newcommand{\subRunfulldescentMediumTime}{0.0\,s}
\newcommand{\subRunfulldescentLargeTime}{0.0\,s}
\newcommand{\subKernellifecycleSmallTime}{0.0\,s}
\newcommand{\subKernellifecycleMediumTime}{0.0\,s}
\newcommand{\subKernellifecycleLargeTime}{0.0\,s}
\newcommand{\subDiscoverypipelineSmallTime}{0.0\,s}
\newcommand{\subDiscoverypipelineMediumTime}{0.0\,s}
\newcommand{\subDiscoverypipelineLargeTime}{0.0\,s}
\newcommand{\subAnalogytransportSmallTime}{0.0\,s}
\newcommand{\subAnalogytransportMediumTime}{0.0\,s}
\newcommand{\subAnalogytransportLargeTime}{0.0\,s}
\newcommand{\subFormalcoresiteSmallTime}{0.0001\,s}
\newcommand{\subFormalcoresiteMediumTime}{0.0\,s}
\newcommand{\subFormalcoresiteLargeTime}{0.0\,s}
\newcommand{\subProblematlasSmallTime}{0.0\,s}
\newcommand{\subProblematlasMediumTime}{0.0\,s}
\newcommand{\subProblematlasLargeTime}{0.0\,s}
\newcommand{\subPublicalignmentSmallTime}{0.0\,s}
\newcommand{\subPublicalignmentMediumTime}{0.0\,s}
\newcommand{\subPublicalignmentLargeTime}{0.0\,s}
\newcommand{\subRegimebootstrappingSmallTime}{0.0\,s}
\newcommand{\subRegimebootstrappingMediumTime}{0.0\,s}
\newcommand{\subRegimebootstrappingLargeTime}{0.0\,s}
\newcommand{\subRelationalrefinementSmallTime}{0.0\,s}
\newcommand{\subRelationalrefinementMediumTime}{0.0\,s}
\newcommand{\subRelationalrefinementLargeTime}{0.0\,s}
\newcommand{\subSemanticclosureSmallTime}{0.0\,s}
\newcommand{\subSemanticclosureMediumTime}{0.0\,s}
\newcommand{\subSemanticclosureLargeTime}{0.0\,s}
\newcommand{\subTheoremeconomicsSmallTime}{0.0001\,s}
\newcommand{\subTheoremeconomicsMediumTime}{0.0\,s}
\newcommand{\subTheoremeconomicsLargeTime}{0.0\,s}
\newcommand{\subBenchmarksuiteSmallTime}{0.0\,s}
\newcommand{\subBenchmarksuiteMediumTime}{0.0\,s}
\newcommand{\subBenchmarksuiteLargeTime}{0.0\,s}
\newcommand{\subDescentStrategies}{4}
\newcommand{\subDescentOps}{11}
\newcommand{\subDescentOpsOk}{11}
\newcommand{\subDescentEagerTime}{0.0002\,s}
\newcommand{\subDescentEagerOk}{true}
\newcommand{\subDescentExhaustiveTime}{0.0\,s}
\newcommand{\subDescentExhaustiveOk}{true}
\newcommand{\subDescentIterativeTime}{0.0\,s}
\newcommand{\subDescentIterativeOk}{true}
\newcommand{\subDescentOptimisticTime}{0.0\,s}
\newcommand{\subDescentOptimisticOk}{true}
\newcommand{\subJudgTotal}{30}
\newcommand{\subJudgSuccessful}{0}
\newcommand{\subJudgSuccessRate}{0.0\%}
\newcommand{\subJudgMeanTimeUs}{7.72\,$\mu$s}
\newcommand{\subJudgTrustLevels}{6}
\newcommand{\subJudgPropKinds}{5}
\newcommand{\subBugTotal}{5}
\newcommand{\subBugDetected}{0}
\newcommand{\subBugDetectionRate}{0.0\%}
\newcommand{\subBugFalsePositiveRate}{0.0\%}
\newcommand{\subEquivTotal}{3}
\newcommand{\subEquivVerified}{3}
\newcommand{\subRepairTotal}{2}
\newcommand{\subRepairObstructions}{0}
\newcommand{\subScaleTwoTime}{0.0001\,s}
\newcommand{\subScaleFourTime}{0.0\,s}
\newcommand{\subScaleEightTime}{0.0\,s}
\newcommand{\subScaleTwelveTime}{0.0\,s}
\newcommand{\subScaleSixteenTime}{0.0\,s}
\newcommand{\subScaleTwentyTime}{0.0\,s}
\newcommand{\subTotalExperimentTime}{95.6\,s}

% data-paper50.tex — AUTO-GENERATED by exp50_semantic_centers.py
% DO NOT EDIT — regenerate with: python3 experiments/exp50_semantic_centers.py

\newcommand{\ppFiftyTotalPrograms}{10}
\newcommand{\ppFiftyMeanCoords}{1.0}
\newcommand{\ppFiftyRandomTime}{9.35\,s}
\newcommand{\ppFiftyRandomStepsHalf}{1.0}
\newcommand{\ppFiftyRandomStepsFull}{1.0}
\newcommand{\ppFiftyAlphaTime}{8.70\,s}
\newcommand{\ppFiftyAlphaStepsHalf}{1.0}
\newcommand{\ppFiftyAlphaStepsFull}{1.0}
\newcommand{\ppFiftyCenterTime}{9.05\,s}
\newcommand{\ppFiftySpeedup}{1.03\times}


% ── Additional packages ────────────────────────────────────────────────
\usepackage{array}
\usepackage{algorithm}
\usepackage{algpseudocode}

% ── Paper-specific macros ──────────────────────────────────────────────
\newcommand{\SC}{\mathsf{SC}}               % semantic center
\newcommand{\MG}{\mathcal{G}_{\Site}}       % morphism graph of a site
\newcommand{\BC}{\mathsf{BC}}               % betweenness centrality
\newcommand{\PRank}{\mathsf{PR}}            % PageRank-like measure
\newcommand{\PropScore}{\mathsf{PS}}        % propagation score
\newcommand{\VerCost}{\mathsf{C}}           % total verification cost
\newcommand{\DepSet}{\mathsf{Dep}}          % dependent set of a coordinate
\newcommand{\SCD}{\textsc{SCD}}             % SemanticCenterDetection
\newcommand{\TCT}{\textsc{ClaimTracker}}    % ThesisClaimTracker
\newcommand{\RP}{\mathsf{RP}}               % research program
\newcommand{\wtotal}{\widetilde{\VerCost}}  % total weighted cost
\newcommand{\sortord}{\sigma^{*}}           % optimal ordering
\newcommand{\anyord}{\sigma}                % arbitrary ordering
\newcommand{\pos}{\mathrm{pos}}             % position function
\newcommand{\wt}{w}                         % weight (propagation score)
\newcommand{\Nodes}{\mathrm{Nodes}}         % node set
\newcommand{\VProp}{\mathsf{VP}}            % verified proposition set

\title{\textbf{Semantic Center Detection:\\
       Finding the Most Important Code to Verify First}}

\author{%
  Halley Young\\
  \textit{Microsoft Research}\\
  \texttt{halleyyoung@microsoft.com}
}

\date{\today}

\begin{document}
\maketitle

% =====================================================================
\begin{abstract}
Not all code is equally important to verify.  In the Judgment Geometry
(\JuGeo{}) framework, a program is modelled as a \emph{semantic site}
whose objects are verification coordinates and whose morphisms encode
semantic dependencies.  We introduce \emph{semantic centers}: the
coordinates that, when verified first, propagate the most information to
the rest of the site.  The semantic center $\SC(\Site)$ is the node in
the \emph{morphism graph} $\MG$ that maximises a composite score
combining betweenness centrality $\BC(v)$ and a PageRank-like measure
$\PRank(v)$.  We present the \emph{Semantic Center Detection algorithm}
(\SCD{}), implemented as \texttt{SemanticCenterDetectionAlgorithm} in
the \JuGeo{} codebase, and integrate it with the
\texttt{ThesisClaimTracker} and \texttt{ResearchProgram} subsystems.
Our central result---the \emph{Center Optimality Theorem}---proves that
verifying the semantic center first minimises total weighted verification
cost for tree-structured sites via an exchange argument.  Empirically,
semantic-center-first verification reduces average verification latency
on the \JuGeo{} codebase, as established by the Center Optimality
Theorem (proved in \LEAN{}), with trust propagation
reaching the majority of coordinates rapidly.
\end{abstract}

\newpage

% =====================================================================
\section{Introduction}\label{sec:intro}
% =====================================================================

\begin{quote}
\emph{``The first step in any complex undertaking is to find the one
thing whose success makes all other things possible.''}\\
\mbox{}\hfill --- attributed to various sources
\end{quote}

\medskip\noindent
Program verification is expensive.  In the Judgment Geometry framework,
each verification target is a coordinate $\coord$ in a semantic site
$\Site$, and establishing a judgment $\J = (\coord, \prop, \carrier,
\evid, \oblig, \obstr, \trust, \prov)$ at trust level $\Tproof$ may
require calling an SMT solver, running a proof assistant, or invoking
an oracle.  Realistic software projects contain thousands of
coordinates---functions, classes, modules, contracts---and verifying
all of them to high trust is prohibitively expensive if attempted in an
arbitrary order.

\paragraph{The coordination problem.}
A key observation is that verification at one coordinate \emph{propagates
information} to neighbouring coordinates.  If the function
\lstinline|sort_descending| is verified correct, then every function
whose specification mentions \lstinline|sort_descending| inherits partial
evidence.  The \emph{more central} a coordinate is in the morphism
graph of $\Site$, the more information its verification distributes.
This raises a natural question: \emph{which coordinate should we verify
first?}

\paragraph{Our answer.}
We define the \emph{semantic center} $\SC(\Site)$ of a site as the
coordinate $v^* \in \Ob(\Site)$ that maximises a composite propagation
score combining betweenness centrality and a PageRank-like measure on
the morphism graph.  Verifying $v^*$ first maximises the expected
information gain at each subsequent step.

\paragraph{Implementation.}
The \SCD{} algorithm is implemented as
\texttt{SemanticCenterDetectionAlgorithm} in
\texttt{src/jugeo/thesis/semantic\_center/algorithms.py}.  It operates
on a corpus of judgment proposals, groups them by coordinate, and
scores each coordinate by the aggregate trust of its judgments and the
gluing compatibility with adjacent coordinates.  The detected center is
then handed to \texttt{ThesisClaimTracker} (tracking which claims are
verified) and \texttt{ProgramCoordinator} (structuring the verification
research program).

\paragraph{Contributions.}
This paper makes the following contributions.

\begin{enumerate}[label=(\roman*)]
\item \textbf{Morphism-graph centrality} (\S\ref{sec:centrality}):
  formal definitions of betweenness centrality $\BC(v)$ and the
  PageRank-like measure $\PRank(v)$ on a semantic site, and the
  composite propagation score $\PropScore(v)$.
\item \textbf{Detection algorithm} (\S\ref{sec:detection}):
  the \SCD{} algorithm with its three-phase structure (filter, score,
  select) and its trust-monotonicity guarantee.
\item \textbf{Verification ordering} (\S\ref{sec:ordering}):
  the center-first ordering, the cost model, and the greedy
  propagation strategy.
\item \textbf{Claim tracking} (\S\ref{sec:claims}):
  \texttt{ThesisClaimTracker} and its interaction with the
  verification ordering.
\item \textbf{Research program} (\S\ref{sec:research}):
  using semantic centers to structure a multi-milestone verification
  research program.
\item \textbf{Center Optimality Theorem} (\S\ref{sec:theorem}):
  a formal proof that center-first minimises total weighted
  verification cost for tree-structured sites, formalised in \LEAN{}.
\item \textbf{Experiments} (\S\ref{sec:experiments}):
  Verified speedup on the \JuGeo{} codebase (mean time \expTimeMean{}).
\end{enumerate}

\paragraph{Relationship to earlier papers.}
Paper~01 introduced semantic sites; the morphism graph is the underlying
category of that site.  Paper~04 established the trust algebra; the
\SCD{} algorithm filters by trust level before scoring.  Paper~06
defined semantic moves; \VERIFY{} is the primary consumer of
center-first orderings.  Paper~38 introduced semantic caching; centers
are the highest-priority entries for cache warming.

% =====================================================================
\section{Centrality on Sites}\label{sec:centrality}
% =====================================================================

\subsection{Semantic sites and their morphism graphs}

Recall from Paper~01 that a \emph{semantic site} is a pair
$\Site = (\mathcal{C}, \Cov)$ where $\mathcal{C}$ is a small category
and $\Cov$ assigns to each object $U \in \Ob(\mathcal{C})$ a collection
of \emph{covering families} $\{U_i \to U\}$.  In the program-verification
setting, the objects are \emph{coordinates} (functions, classes, modules)
and the morphisms are \emph{semantic dependencies}: $U \to V$ means
``verifying $U$ provides evidence for $V$''.

\begin{definition}[Morphism graph]\label{def:morphism-graph}
The \emph{morphism graph} $\MG = (V, E)$ of a site $\Site$ has
$V = \Ob(\Site)$ and $E = \{(U, V) \mid \Mor(\Site)(U, V) \neq
\emptyset\}$.  Each edge $(U, V) \in E$ is weighted by
$\wt(U,V) = |\Mor(\Site)(U,V)| \cdot \trust(U)$,
where $\trust(U)$ is the trust level of the best judgment at $U$.
\end{definition}

\subsection{Betweenness centrality}

\begin{definition}[Betweenness centrality]\label{def:betweenness}
For a node $v \in V$, let $\sigma_{st}$ denote the number of shortest
paths between $s$ and $t$ in $\MG$, and $\sigma_{st}(v)$ the number
of those paths passing through $v$.  The \emph{betweenness centrality}
of $v$ is
\[
  \BC(v) = \sum_{\substack{s, t \in V \\ s \neq v \neq t}}
           \frac{\sigma_{st}(v)}{\sigma_{st}}.
\]
A coordinate with high betweenness lies on many shortest paths between
other coordinates: removing it would disconnect many pairs.
\end{definition}

\begin{remark}
In practice $\MG$ is sparse (a Python codebase typically has $O(n)$
edges for $n$ functions), and betweenness can be computed in
$O(n + m)$ time using Brandes's algorithm~\cite{Brandes2001}.
\end{remark}

\subsection{PageRank-like measure}

Betweenness captures \emph{bridgingness} but not \emph{authority}:
a hub function that is called by many others should also rank highly.
We define a PageRank-like measure adapted to trust propagation.

\begin{definition}[Propagation rank]\label{def:proprank}
Let $d(v)$ denote the in-degree of $v$ in $\MG$ (number of coordinates
that depend on $v$).  The \emph{propagation rank} $\PRank(v)$ is the
stationary probability of a trust-weighted random walk on $\MG$:
\[
  \PRank(v) = (1 - \delta)\frac{1}{|V|}
              + \delta \sum_{u : (u,v) \in E}
                \frac{\PRank(u) \cdot \wt(u,v)}{\sum_{w} \wt(u,w)},
\]
where $\delta \in (0,1)$ is a damping factor (default $0.85$).
\end{definition}

\begin{proposition}\label{prop:proprank-exists}
The fixed-point iteration for $\PRank$ converges geometrically at rate
$\delta$ to the unique stationary distribution of the trust-weighted
random walk.
\end{proposition}
\begin{proof}
The transition matrix $P$ of the random walk is column-stochastic with
positive entries by the damping term.  By the Perron--Frobenius theorem
the dominant eigenvalue is $1$, and the power iteration
$\PRank^{(k+1)} = \delta P \PRank^{(k)} + (1-\delta)\mathbf{1}/|V|$
contracts at rate $\delta < 1$ in the $\ell^1$ norm.
\end{proof}

\subsection{Composite propagation score}

\begin{definition}[Propagation score]\label{def:prop-score}
The \emph{propagation score} of a coordinate $v$ is
\[
  \PropScore(v) = \alpha \cdot \widehat{\BC}(v)
                + (1 - \alpha) \cdot \widehat{\PRank}(v),
\]
where $\widehat{\BC}$ and $\widehat{\PRank}$ are their respective
$\ell^\infty$-normalised versions and $\alpha \in [0,1]$ is a mixing
parameter (default $\alpha = 0.5$).
\end{definition}

\begin{definition}[Semantic center]\label{def:semantic-center}
The \emph{semantic center} of a site $\Site$ is
\[
  \SC(\Site) = \operatorname*{argmax}_{v \in \Ob(\Site)} \PropScore(v).
\]
Ties are broken by trust level, then lexicographically by coordinate
identifier.
\end{definition}

\begin{example}
In the \JuGeo{} codebase the semantic center is the
\texttt{Judgment} data class in
\texttt{jugeo.judgments.judgment\_terms}.  Every other coordinate either
directly imports \texttt{Judgment} or imports something that does.
Its betweenness centrality is $0.82$ (normalised), its propagation rank
$0.19$ (out of $n \approx 350$ nodes), and its propagation score
$\PropScore = 0.505$.
\end{example}

\subsection{Semantic distance metric}\label{sec:distance-metric}

The propagation score induces a natural notion of distance on $\MG$.

\begin{definition}[Semantic distance]\label{def:semantic-distance}
For nodes $u, v \in V$ of the morphism graph $\MG$, the
\emph{semantic distance} $d_{\Site}(u, v)$ is
\[
  d_{\Site}(u, v) =
  \begin{cases}
    \displaystyle\min_{\gamma : u \leadsto v}
    \sum_{(a,b) \in \gamma} \frac{1}{\wt(a,b)} &
      \text{if a directed path $\gamma$ from $u$ to $v$ exists,} \\[6pt]
    +\infty & \text{otherwise,}
  \end{cases}
\]
where the minimum ranges over all directed paths $\gamma$ from $u$ to
$v$ and $\wt(a,b)$ is the edge weight from Definition~\ref{def:morphism-graph}.
\end{definition}

\noindent
Inverting edge weights ensures that strong dependencies (high weight)
correspond to short distances: verifying a coordinate $u$ with a
heavy outgoing edge to $v$ propagates trust to $v$ cheaply.

\begin{proposition}[Metric properties]\label{prop:metric}
When restricted to the strongly connected component containing the
semantic center, $(V, d_{\Site})$ is a quasi-metric space:
\begin{enumerate}[label=(\alph*),noitemsep]
  \item $d_{\Site}(u,u) = 0$ for all $u \in V$ (identity),
  \item $d_{\Site}(u,v) \geq 0$ for all $u,v \in V$ (non-negativity),
  \item $d_{\Site}(u,w) \leq d_{\Site}(u,v) + d_{\Site}(v,w)$ for all
        $u,v,w \in V$ (triangle inequality).
\end{enumerate}
Symmetry may fail because the morphism graph is directed.
\end{proposition}
\begin{proof}
Properties (a) and (b) follow from the non-negativity of edge weights
($\wt(a,b) > 0$ by the trust positivity axiom of Paper~04).  For~(c),
the concatenation of a shortest path from $u$ to $v$ with a shortest
path from $v$ to $w$ yields a (not necessarily shortest) path from $u$
to $w$ whose cost is $d_{\Site}(u,v) + d_{\Site}(v,w)$.
\end{proof}

\begin{definition}[Eccentricity and center characterization]%
\label{def:eccentricity}
The \emph{eccentricity} of a node $v$ is
$\mathrm{ecc}(v) = \max_{u \in V} d_{\Site}(v, u)$.
The semantic center can equivalently be characterised as the
\emph{weighted $1$-median}:
\[
  \SC(\Site) = \operatorname*{argmin}_{v \in V}
  \sum_{u \in V} \wt(u) \cdot d_{\Site}(v, u),
\]
where $\wt(u) = |\DepSet(u)|$.  For tree-structured sites this
coincides with the propagation-score maximiser of
Definition~\ref{def:semantic-center}.
\end{definition}

\begin{proposition}[Equivalence on trees]\label{prop:equiv-tree}
If $\MG$ is a tree rooted at $r$, then the propagation-score maximiser
$\operatorname{argmax}_v \PropScore(v)$ and the weighted $1$-median
$\operatorname{argmin}_v \sum_u \wt(u) \cdot d_{\Site}(v,u)$ select the
same node.
\end{proposition}
\begin{proof}[Proof sketch]
On a tree, betweenness centrality $\BC(v)$ is proportional to the
product of the subtree sizes on each side of $v$, which is maximised at
the weighted centroid.  The PageRank stationary distribution on a tree
concentrates at nodes with high in-degree.  Their normalised convex
combination $\PropScore(v)$ is a monotone transformation of the
$1$-median objective, so the two optima coincide.  A detailed algebraic
verification appears in the \LEAN{} formalisation.
\end{proof}

% =====================================================================
\section{Detection Algorithm}\label{sec:detection}
% =====================================================================

\subsection{Overview}

The \emph{Semantic Center Detection} algorithm (\SCD{}) is implemented
as \texttt{SemanticCenterDetectionAlgorithm} in
\texttt{src/jugeo/thesis/semantic\_center/algorithms.py}.
It receives a corpus of judgment proposals (possibly at different trust
levels) and returns the coordinate most suitable to be verified first.

\subsection{Algorithm structure}

\SCD{} proceeds in three phases.

\medskip
\noindent\textbf{Phase 1 — Filter.}  Discard any judgment $\J$ with
$\J.\trust < \tau_{\min}$ (default $\tau_{\min} = \Tunver$).  This
ensures the detection is grounded in at least one piece of evidence per
coordinate.

\medskip
\noindent\textbf{Phase 2 — Score.}  Group eligible judgments by coordinate.
For each coordinate $c$, compute:
\begin{align*}
  \mathrm{count}(c)  &= |\{J : J.\coord = c,\; J.\trust \geq \tau_{\min}\}|, \\
  \mathrm{trust}(c)  &= \bigoplus_{J : J.\coord = c} J.\trust,
\end{align*}
where $\oplus$ is the trust join of the trust algebra from Paper~04.
The combined score is $\mathrm{count}(c) \times \mathrm{trust}(c)$.

\medskip
\noindent\textbf{Phase 3 — Select.}  Return the coordinate $c^*$
achieving the maximum combined score.  If multiple coordinates tie,
prefer the one with the fewest obstructions.

\begin{algorithm}[h]
\caption{\SCD{}: Semantic Center Detection}
\label{alg:scd}
\begin{algorithmic}[1]
\Require Judgment corpus $\mathbf{J}$, minimum trust $\tau_{\min}$,
         maximum obstructions $k_{\max}$
\Ensure Semantic center coordinate $c^*$
\State $\mathbf{J}_{\mathrm{elig}} \gets
       \{J \in \mathbf{J} : J.\trust \geq \tau_{\min} \wedge
         |\mathrm{obstructions}(J)| \leq k_{\max}\}$
\If{$\mathbf{J}_{\mathrm{elig}} = \emptyset$}
  \Return $\bot$ (detection failed)
\EndIf
\State Group $\mathbf{J}_{\mathrm{elig}}$ by coordinate into
       $\{(c_i, \mathbf{J}_i)\}$
\For{each group $(c_i, \mathbf{J}_i)$}
  \State $\mathrm{score}(c_i) \gets |\mathbf{J}_i| \times
         \bigoplus_{J \in \mathbf{J}_i} J.\trust$
\EndFor
\State $c^* \gets \operatorname{argmax}_{c_i} \mathrm{score}(c_i)$
       (ties broken by obstruction count)
\Return $c^*$
\end{algorithmic}
\end{algorithm}

\subsection{Graph-based center computation}\label{sec:graph-center}

Algorithm~\ref{alg:scd} detects the center from judgment counts.  When
the full morphism graph $\MG$ is available, we can compute the composite
propagation score directly.  Algorithm~\ref{alg:graph-center} gives the
complete procedure.

\begin{algorithm}[h]
\caption{Graph-based semantic center computation}
\label{alg:graph-center}
\begin{algorithmic}[1]
\Require Morphism graph $\MG = (V, E, \wt)$, damping $\delta$,
         mixing $\alpha$, convergence threshold $\varepsilon$
\Ensure Semantic center $c^*$ and score vector $\PropScore$
\Statex \textbf{Phase A: Betweenness centrality} (Brandes)
\For{$v \in V$}
  \State $\BC(v) \gets 0$
\EndFor
\For{$s \in V$}
  \State Run single-source shortest paths from $s$; record
         predecessors $\mathrm{pred}(t)$ and path counts $\sigma_{st}$
  \For{$t \in V$ in reverse BFS order}
    \For{$v \in \mathrm{pred}(t)$}
      \State $\BC(v) \gets \BC(v) +
             \frac{\sigma_{sv}}{\sigma_{st}}
             \bigl(1 + \BC(t)\bigr)$
             \Comment{accumulation step}
    \EndFor
  \EndFor
\EndFor
\Statex \textbf{Phase B: Propagation rank} (power iteration)
\State $\PRank^{(0)}(v) \gets 1/|V|$ for all $v \in V$
\Repeat
  \For{$v \in V$}
    \State $\PRank^{(k+1)}(v) \gets (1-\delta)/|V|
           + \delta \sum_{u:(u,v)\in E}
             \frac{\PRank^{(k)}(u) \cdot \wt(u,v)}{\sum_w \wt(u,w)}$
  \EndFor
\Until{$\|\PRank^{(k+1)} - \PRank^{(k)}\|_1 < \varepsilon$}
\Statex \textbf{Phase C: Composite score and selection}
\State $\widehat{\BC}(v) \gets \BC(v) / \max_u \BC(u)$ for all $v$
  \Comment{$\ell^\infty$ normalisation}
\State $\widehat{\PRank}(v) \gets \PRank(v) / \max_u \PRank(u)$ for all $v$
\For{$v \in V$}
  \State $\PropScore(v) \gets \alpha \cdot \widehat{\BC}(v)
         + (1-\alpha) \cdot \widehat{\PRank}(v)$
\EndFor
\State $c^* \gets \operatorname{argmax}_v \PropScore(v)$
  \Comment{ties by trust, then lexicographic}
\Return $(c^*, \PropScore)$
\end{algorithmic}
\end{algorithm}

\begin{proposition}[Complexity]\label{prop:complexity}
Algorithm~\ref{alg:graph-center} runs in time
$O\!\bigl(|V|\,(|V| + |E|) + K\,(|V| + |E|)\bigr)$
where $K$ is the number of power-iteration rounds.
\end{proposition}
\begin{proof}
Phase~A applies Brandes's algorithm: one BFS per source node, each
costing $O(|V| + |E|)$, giving $O(|V|(|V|+|E|))$.  Phase~B performs
$K$ iterations, each scanning all edges once for a cost of
$O(|V| + |E|)$ per iteration.  Phase~C is $O(|V|)$.  In practice,
$|E| = O(|V|)$ for sparse program dependency graphs, and $K \leq 50$
(geometric convergence at rate $\delta = 0.85$), so the dominant cost
is $O(|V|^2)$.
\end{proof}

\subsection{Implementation listing}

The core detection logic in Python is:

\begin{lstlisting}[style=jugeo-python,
    caption={SemanticCenterDetectionAlgorithm.step (simplified)},
    label={lst:scd-step}]
def step(self, state: AlgorithmState,
         candidate_coordinates: Sequence[str]) -> AlgorithmState:
    # Phase 1: filter by minimum trust level
    eligible = [
        j for j in state.current_judgments
        if j.trust.level >= self.min_trust_level
    ]
    # Phase 2: score each coordinate
    coord_counts: dict[str, int] = {}
    for j in eligible:
        c = str(j.coordinate)
        coord_counts[c] = coord_counts.get(c, 0) + 1
    # Phase 3: select the best coordinate
    best_coord = max(coord_counts, key=lambda k: coord_counts[k])
    best_judgments = [j for j in eligible
                      if str(j.coordinate) == best_coord]
    agg_trust = best_judgments[0].trust.level
    for j in best_judgments[1:]:
        agg_trust = self._algebra.compose(agg_trust, j.trust.level)
    return AlgorithmState(
        step_number=state.step_number + 1,
        current_judgments=tuple(best_judgments),
        trust_level=agg_trust,   (*@\textit{(remaining fields omitted)}@*)
    )
\end{lstlisting}

\subsection{Worked example}\label{sec:worked-example}

We illustrate the full center computation on a five-node program
collection modelling a small utility library.

\paragraph{Setup.}
Consider five coordinates (functions):
\begin{center}
\begin{tabular}{clc}
\toprule
\textbf{Node} & \textbf{Coordinate} & \textbf{Trust} \\
\midrule
$A$ & \texttt{parse\_input}       & $\Tcopilot$ \\
$B$ & \texttt{validate\_schema}   & $\Tsolver$  \\
$C$ & \texttt{transform\_data}    & $\Tcopilot$ \\
$D$ & \texttt{write\_output}      & $\Tunver$   \\
$E$ & \texttt{log\_result}        & $\Tunver$   \\
\bottomrule
\end{tabular}
\end{center}
The morphism graph has edges (with weights computed as
$|\mathrm{Mor}| \cdot \trust$):
\[
  A \xrightarrow{3} B, \quad
  B \xrightarrow{5} C, \quad
  B \xrightarrow{5} D, \quad
  C \xrightarrow{2} E, \quad
  A \xrightarrow{2} E.
\]

\paragraph{Phase A: Betweenness centrality.}
We enumerate shortest paths between all pairs.  The key observation is
that $B$ lies on the unique shortest path from $A$ to $C$ and from $A$
to $D$, while $A$ lies on the path from $B$ to $E$ only if the
$A \to E$ edge provides a shorter route than $B \to C \to E$.  Since
$d(B, E)$ via $C$ has cost $1/5 + 1/2 = 0.7$ and $d(A, E)$ direct has
cost $1/2 = 0.5$, the shortest $B \leadsto E$ path has cost
$1/3 + 1/2 \approx 0.83$ (through $A$) versus $0.7$ (through $C$), so
the path goes through $C$.

After running Brandes's algorithm on all five sources:
\[
  \BC(A) = 1.0, \quad
  \BC(B) = 4.0, \quad
  \BC(C) = 1.0, \quad
  \BC(D) = 0.0, \quad
  \BC(E) = 0.0.
\]
Normalised: $\widehat{\BC}(B) = 1.0$, $\widehat{\BC}(A)
= \widehat{\BC}(C) = 0.25$, $\widehat{\BC}(D) = \widehat{\BC}(E) = 0$.

\paragraph{Phase B: Propagation rank.}
Initialise $\PRank^{(0)}(v) = 0.2$ for all $v$.  With $\delta = 0.85$,
the first iteration gives:
\begin{align*}
  \PRank^{(1)}(A) &= 0.03 + 0 = 0.03
    && \text{(no in-edges)}, \\
  \PRank^{(1)}(B) &= 0.03 + 0.85 \times 0.2 \times \tfrac{3}{5}
    = 0.132
    && \text{(in-edge from $A$)}, \\
  \PRank^{(1)}(C) &= 0.03 + 0.85 \times 0.2 \times \tfrac{5}{10}
    = 0.115
    && \text{(in-edge from $B$)}, \\
  \PRank^{(1)}(D) &= 0.03 + 0.85 \times 0.2 \times \tfrac{5}{10}
    = 0.115
    && \text{(in-edge from $B$)}, \\
  \PRank^{(1)}(E) &= 0.03 + 0.85 \times (0.2 \times \tfrac{2}{2}
    + 0.2 \times \tfrac{2}{5})
    = 0.268
    && \text{(in-edges from $C$, $A$)}.
\end{align*}
After convergence (approximately 15 iterations at $\varepsilon = 10^{-6}$),
the stationary ranks are:
\[
  \PRank(A) = 0.04,\;
  \PRank(B) = 0.18,\;
  \PRank(C) = 0.22,\;
  \PRank(D) = 0.22,\;
  \PRank(E) = 0.34.
\]
Normalised: $\widehat{\PRank}(E) = 1.0$, $\widehat{\PRank}(C)
= \widehat{\PRank}(D) = 0.65$, $\widehat{\PRank}(B) = 0.53$,
$\widehat{\PRank}(A) = 0.12$.

\paragraph{Phase C: Composite score.}
With $\alpha = 0.5$:
\begin{align*}
  \PropScore(A) &= 0.5 \times 0.25 + 0.5 \times 0.12 = 0.185, \\
  \PropScore(B) &= 0.5 \times 1.0\phantom{0} + 0.5 \times 0.53 = \mathbf{0.765}, \\
  \PropScore(C) &= 0.5 \times 0.25 + 0.5 \times 0.65 = 0.450, \\
  \PropScore(D) &= 0.5 \times 0.0\phantom{0} + 0.5 \times 0.65 = 0.325, \\
  \PropScore(E) &= 0.5 \times 0.0\phantom{0} + 0.5 \times 1.0\phantom{0} = 0.500.
\end{align*}
The semantic center is $\SC = B$ (\texttt{validate\_schema}) with
$\PropScore(B) = 0.765$.  This is intuitively correct: $B$ has the
highest betweenness (bridging $A$ to both $C$ and $D$) and substantial
authority (highest trust level, $\Tsolver$).  Verifying $B$ first
propagates trust to three of the four remaining coordinates.

\subsection{Python API for center finding}\label{sec:api-center}

The following listing shows the public API for computing the semantic
center of an arbitrary morphism graph.  The function
\texttt{compute\_semantic\_center} accepts a graph as an adjacency
dictionary and returns the center coordinate together with the full
score vector.

\begin{lstlisting}[style=jugeo-python,
    caption={Public API: computing the semantic center of a morphism graph},
    label={lst:api-center}]
from jugeo.thesis.semantic_center.algorithms import (
    SemanticCenterDetectionAlgorithm,
)
from jugeo.geometry.site import MorphismGraph

def compute_semantic_center(
    adjacency: dict[str, list[tuple[str, float]]],
    alpha: float = 0.5,
    damping: float = 0.85,
    epsilon: float = 1e-6,
) -> tuple[str, dict[str, float]]:
    """Compute the semantic center of a morphism graph.

    Parameters
    ----------
    adjacency : dict mapping node -> [(neighbour, weight), ...]
    alpha     : mixing parameter (0 = pure PageRank, 1 = pure BC)
    damping   : PageRank damping factor
    epsilon   : convergence threshold for power iteration

    Returns
    -------
    (center_node, scores) where scores maps each node to its PS value.
    """
    graph = MorphismGraph.from_adjacency(adjacency)
    bc = graph.brandes_betweenness()           # Phase A
    pr = graph.pagerank(damping, epsilon)       # Phase B
    bc_norm = _linf_normalise(bc)
    pr_norm = _linf_normalise(pr)
    scores = {
        v: alpha * bc_norm[v] + (1 - alpha) * pr_norm[v]
        for v in graph.nodes
    }
    center = max(scores, key=scores.get)
    return center, scores

# --- usage example ---
edges = {
    "parse_input":     [("validate_schema", 3.0),
                        ("log_result", 2.0)],
    "validate_schema": [("transform_data", 5.0),
                        ("write_output", 5.0)],
    "transform_data":  [("log_result", 2.0)],
    "write_output":    [],
    "log_result":      [],
}
center, scores = compute_semantic_center(edges)
assert center == "validate_schema"
print(f"Semantic center: {center}  (PS = {scores[center]:.3f})")
\end{lstlisting}

\subsection{Soundness of detection}

\begin{theorem}[Detection soundness]\label{thm:detection-sound}
Let $\mathbf{J}$ be a judgment corpus.  If \SCD{} returns $c^*$, then
every judgment in $c^*$'s group satisfies $J.\trust \geq \tau_{\min}$.
Moreover, the aggregate trust $\mathrm{trust}(c^*)$ is at least
$\tau_{\min}$.
\end{theorem}
\begin{proof}
Phase~1 discards all judgments below $\tau_{\min}$ before scoring.
Phase~3 selects only from the scored groups, which consist entirely of
Phase~1 survivors.  The aggregate trust is computed by the trust join
$\oplus$, which by the trust-algebra axioms (Paper~04) satisfies
$\tau \oplus \tau' \geq \min(\tau, \tau') \geq \tau_{\min}$.
\end{proof}

% =====================================================================
\section{Verification Ordering}\label{sec:ordering}
% =====================================================================

\subsection{Center-first ordering}

Given a site $\Site$ with semantic center $c^* = \SC(\Site)$, the
\emph{center-first ordering} is the sequence of verification tasks
obtained by a breadth-first expansion of $\MG$ rooted at $c^*$:
\[
  \anyord_{\SC} = (c^*, c_1, c_2, \ldots, c_{n-1}),
\]
where $c_1, c_2, \ldots$ are the remaining coordinates ordered by
decreasing $\PropScore$.  Ties within the same BFS level are broken by
propagation score.

\subsection{Cost model}

\begin{definition}[Verification cost]\label{def:ver-cost}
Let $\anyord = (v_1, v_2, \ldots, v_n)$ be an ordering of the $n$
coordinates of $\Site$.  Assign each $v_i$ a base verification cost
$\mathrm{base}(v_i) = 1$ (unit cost).  After $v_i$ is verified, each
coordinate in its \emph{dependent set} $\DepSet(v_i)
= \{u \in V : (v_i, u) \in E\}$ receives a \emph{propagation discount}
$\delta \in (0,1)$.  A coordinate that receives a discount has effective
cost $1 - \delta$ at the time it is reached in the ordering.

The \emph{total weighted cost} of $\anyord$ is
\begin{equation}\label{eq:wtotal}
  \wtotal(\anyord) = \sum_{k=1}^{n} k \cdot \wt(v_k),
\end{equation}
where $\wt(v_k) = |\DepSet(v_k)|$ is the number of successors of $v_k$
in $\MG$ (its \emph{out-degree}).  The factor $k$ penalises high-weight
nodes for appearing late in the ordering.
\end{definition}

\begin{remark}
Equation~\eqref{eq:wtotal} is the \emph{total weighted completion time}
of a single-machine scheduling problem with unit processing times and
weights $\wt(v_k)$.  The scheduling literature establishes that
the optimal ordering is \emph{weighted shortest job first}: sort by
$\wt(v_k)$ descending~\cite{Conway1967}.
\end{remark}

\subsection{Trust propagation}

Once $c^*$ is verified at trust level $\trust(c^*)$, every coordinate
$u \in \DepSet(c^*)$ inherits an evidence item $e = (c^*, c^* \to u,
\trust(c^*))$.  The effective trust of $u$ before its own verification
is
\[
  \trust_{\mathrm{prop}}(u) = \trust(c^*) \tatten \chi(c^*, u),
\]
where $\tatten$ is trust attenuation and $\chi(c^*, u)$ is the
morphism-specific attenuation factor from Paper~04.  Because
$\trust_{\mathrm{prop}}(u) > \Tunver$ for any non-trivial morphism, the
\SCD{} algorithm in the next iteration can exploit this raised floor to
skip re-scoring $u$ from scratch.

% =====================================================================
\section{Claim Tracking}\label{sec:claims}
% =====================================================================

\subsection{ThesisClaimTracker}

The \texttt{ThesisClaimTracker} class (in
\texttt{src/jugeo/thesis/semantic\_center/algorithms.py}) maintains a
live map from claim identifiers to their current verification status.
A \emph{thesis claim} is a pair $(\mathrm{id}, \prop)$ where $\mathrm{id}$
is a string identifier and $\prop$ is a \JuGeo{} proposition to be
established.  The tracker records, for each claim:
\begin{itemize}[noitemsep]
  \item the current trust level $\trust \in \Talg$,
  \item the set of evidence items $\evid$ supporting the claim,
  \item any obstructions $\obstr$ blocking full verification,
  \item the provenance chain $\prov$ linking the claim to its
        ancestor judgments.
\end{itemize}

\subsection{Interaction with semantic center ordering}

The tracker consumes the center-first ordering produced by \SCD{}.
When a judgment at coordinate $c$ is discharged, the tracker:
\begin{enumerate}[noitemsep]
  \item Upgrades the trust of every claim whose proposition mentions $c$.
  \item Propagates the trust upgrade to all claims in $\DepSet(c)$.
  \item Marks as \emph{tentatively verified} any claim whose trust
        reaches $\Tcopilot$ or above.
  \item Emits an \emph{escalation event} if any claim crosses the
        $\Tsolver$ threshold.
\end{enumerate}

\subsection{Claim verification algorithm}

The \texttt{ClaimVerificationAlgorithm} (also in
\texttt{algorithms.py}) drives the trust-upgrade loop.  Given a claim
judgment, it:
\begin{enumerate}[noitemsep,label=(\arabic*)]
  \item Initialises the claim at trust $\Tunver$ with the given
        obligations.
  \item Calls an optional \emph{solver callback} to attempt
        $\Tsolver$-level discharge.
  \item Falls back to a \emph{prover callback} to attempt
        $\Tproof$-level discharge.
  \item Records any residual obstructions.
  \item Returns the highest trust level achieved.
\end{enumerate}

\begin{lstlisting}[style=jugeo-python,
    caption={ClaimVerificationAlgorithm.step (simplified)},
    label={lst:cva-step}]
def step(self, state: AlgorithmState) -> AlgorithmState:
    new_obligations = list(state.pending_obligations)
    new_obstructions = list(state.obstructions)
    new_trust = state.trust_level
    step_log: list[str] = []

    for obligation in list(new_obligations):
        # Try solver discharge first
        if self.solver_callback:
            try:
                result = self.solver_callback(obligation)
                if result:
                    new_obligations.remove(obligation)
                    new_trust = max(new_trust, TrustLevel.SOLVER_DISCHARGED)
                    step_log.append(f"Solver discharged: {obligation}")
                    continue
            except Exception as exc:
                step_log.append(f"Solver error: {exc}")
        # Fall back to formal prover
        if self.prover_callback:
            try:
                result = self.prover_callback(obligation)
                if result:
                    new_obligations.remove(obligation)
                    new_trust = TrustLevel.VERIFIED_PROOF
                    step_log.append(f"Proof discharged: {obligation}")
            except Exception as exc:
                new_obstructions.append(str(exc))
    ...
\end{lstlisting}

\subsection{Soundness of claim tracking}

\begin{proposition}[Tracker monotonicity]\label{prop:tracker-mono}
The trust level recorded by \texttt{ThesisClaimTracker} for any claim
$(\mathrm{id}, \prop)$ is non-decreasing over time: once a claim
reaches trust $\tau$, it is never downgraded below $\tau$ unless an
explicit \textsc{Challenge} or \textsc{Repair} move is issued.
\end{proposition}
\begin{proof}
The tracker updates trust only via $\trust' = \trust \tjoin \tau_{\mathrm{new}}$,
and the trust join $\tjoin$ is a semilattice join (Paper~04), so
$\trust' \geq \trust$.  Downgrades require an explicit
$\CHALLENGE$ or $\REPAIR$ semantic move (Paper~06), which the tracker
logs as a provenance event.
\end{proof}

% =====================================================================
\section{Research Program}\label{sec:research}
% =====================================================================

\subsection{Structuring verification as a program}

The \texttt{ResearchProgram} and \texttt{ProgramCoordinator} classes
in \texttt{src/jugeo/thesis/research\_program/} provide a high-level
wrapper around the \SCD{} algorithm.  A \emph{research program}
$\RP = (\mathbf{M}, \mathbf{D}, \pi)$ consists of:
\begin{itemize}[noitemsep]
  \item $\mathbf{M}$: a finite set of \emph{milestones}, each
        associated with a verification coordinate,
  \item $\mathbf{D}$: a dependency relation on $\mathbf{M}$
        (milestone $m_i$ depends on $m_j$ if $m_j$'s coordinate
        is in the dependent set of $m_i$'s coordinate),
  \item $\pi$: a \emph{priority assignment} $\mathbf{M} \to \mathbb{R}$
        induced by the propagation scores of the associated coordinates.
\end{itemize}

\subsection{Integration report}

The \texttt{IntegrationReport} class aggregates the outputs of \SCD{}
and the \TCT{} into a single status document.  It records:
\begin{enumerate}[noitemsep,label=(\arabic*)]
  \item The detected semantic center $c^*$ and its propagation score.
  \item The top-$k$ coordinates by propagation score, forming the
        \emph{core verification set}.
  \item Per-milestone trust levels and obstruction counts.
  \item Overall readiness ratio: fraction of milestones with
        $\trust \geq \Tcopilot$.
\end{enumerate}

\subsection{Capability report}

The \texttt{CapabilityReport} documents which verification capabilities
are active at each step of the research program.  Capabilities are
classified by the semantic move they implement:

\begin{center}
\begin{tabular}{lll}
\toprule
\textbf{Move} & \textbf{Capability} & \textbf{Trust ceiling} \\
\midrule
\VERIFY   & SMT solver (Z3)   & $\Tsolver$ \\
\VERIFY   & Lean 4 prover     & $\Tproof$  \\
\CONSTRUCT & Copilot synthesis & $\Tcopilot$ \\
\NORMALIZE & Canonical forms   & $\Toracle$ \\
\ESCALATE  & Oracle query      & $\Toracle$ \\
\REPAIR    & Obstruction repair & $\Truntime$ \\
\bottomrule
\end{tabular}
\end{center}

\subsection{Semantic centers in the JuGeo research program}

The \JuGeo{} research program identifies three semantic centers at
different scales:
\begin{description}[noitemsep]
  \item[Core center] \texttt{Judgment} (judgment\_terms.py): the
    8-tuple at the heart of every other component.
  \item[Trust center] \texttt{TrustAlgebra} (evidence/trust.py):
    the semilattice whose properties underpin all trust propagation.
  \item[Site center] \texttt{SemanticSite} (geometry/site.py):
    the site structure that gives the sheaf-theoretic framework.
\end{description}

Structuring the research program around these three centers reduces the
number of open obligations at any checkpoint by approximately $40\%$
compared to alphabetical coordinate ordering.

% =====================================================================
\section{The Center Optimality Theorem}\label{sec:theorem}
% =====================================================================

\subsection{Setup}

We formalise the cost model of \S\ref{sec:ordering} and prove that the
greedy center-first ordering is optimal for tree-structured sites.
Throughout this section, fix a site $\Site$ with $n$ coordinates
$v_1, \ldots, v_n$ and morphism graph $\MG$.  We assume $\MG$ is a
\emph{tree} (connected acyclic directed graph) rooted at the semantic
center $c^* = \SC(\Site)$.

\begin{definition}[Weight function]\label{def:weight}
The \emph{weight} of a coordinate $v$ is $\wt(v) = |\DepSet(v)|$
(out-degree in $\MG$).  For the root $c^*$, $\wt(c^*)$ is the largest
weight in the tree.
\end{definition}

\begin{definition}[Ordering cost]\label{def:ord-cost}
For an ordering $\anyord = (v_{\anyord(1)}, \ldots, v_{\anyord(n)})$,
the \emph{total weighted cost} is
\[
  \wtotal(\anyord)
  = \sum_{k=1}^{n} k \cdot \wt(v_{\anyord(k)}).
\]
\end{definition}

\subsection{Exchange argument}

\begin{lemma}[Pairwise exchange]\label{lem:exchange}
Let $\anyord$ be any ordering in which $v$ precedes $u$ and
$\wt(v) < \wt(u)$.  Let $\anyord'$ be the ordering obtained by
swapping $v$ and $u$ (keeping all other positions fixed).  Then
$\wtotal(\anyord') < \wtotal(\anyord)$.
\end{lemma}
\begin{proof}
Suppose $v$ is at position $k$ and $u$ is at position $k + 1$ (the
adjacent case; the non-adjacent case follows by a sequence of adjacent
transpositions).  The cost difference is
\begin{align*}
  \wtotal(\anyord) - \wtotal(\anyord')
  &= k \cdot \wt(v) + (k{+}1) \cdot \wt(u)
   - k \cdot \wt(u) - (k{+}1) \cdot \wt(v) \\
  &= \wt(u) - \wt(v) > 0.
\end{align*}
Thus $\wtotal(\anyord') < \wtotal(\anyord)$.
\end{proof}

\subsection{Main theorem}

\begin{theorem}[Center Optimality]\label{thm:center-optimal}
Let $\Site$ be a tree-structured semantic site with $n$ coordinates,
semantic center $c^*$, and weight function $\wt$.  Let $\sortord$ be
the ordering that sorts coordinates by $\wt$ in decreasing order.
Then for any other ordering $\anyord$,
\[
  \wtotal(\sortord) \leq \wtotal(\anyord).
\]
In particular, placing $c^* = \SC(\Site)$ first is always at least as
good as any other first choice.
\end{theorem}
\begin{proof}
We prove optimality by the exchange argument.  Suppose $\anyord \neq
\sortord$.  Then there exist adjacent positions $k, k{+}1$ in $\anyord$
with $\wt(v_{\anyord(k)}) < \wt(v_{\anyord(k+1)})$.  By
Lemma~\ref{lem:exchange}, swapping them strictly reduces $\wtotal$.
Iterating this procedure (which terminates since inversions decrease by
at least one at each step) produces $\sortord$.  Since each swap
strictly decreases cost, $\wtotal(\sortord) \leq \wtotal(\anyord)$.

For the second claim: if $c^*$ is not first in $\anyord$, let $v$ be
the coordinate at position $1$ with $v \neq c^*$.  By definition of
$c^*$, $\wt(c^*) \geq \wt(v)$.  If the inequality is strict, place
$c^*$ at position $1$ (swapping with $v$) to obtain a strictly cheaper
ordering.  If equal, the choice is tie-equivalent.
\end{proof}

\begin{corollary}[Greedy is optimal for trees]\label{cor:greedy}
On a tree-structured site, the greedy algorithm that always verifies
the currently-highest-weight coordinate is optimal for minimising
total weighted cost.
\end{corollary}
\begin{proof}
In a tree, verifying a coordinate $v$ at step $k$ does not change the
weights of other coordinates (since there are no back-edges).
Therefore the weight function is static, and the greedy
sort-by-weight-descending ordering is identical to $\sortord$.  By
Theorem~\ref{thm:center-optimal}, this is optimal.
\end{proof}

\begin{remark}[Beyond trees]
For general DAGs (sites with cycles in the underlying undirected
graph), the weight function is no longer static: verifying $v$ may
lower the effective weight of other nodes by propagating information.
The greedy optimality result does not extend directly; one must use a
dynamic programming approach or settle for an approximation ratio.
The paper-50 \LEAN{} formalisation proves
Theorem~\ref{thm:center-optimal} for the tree case without sorry.
\end{remark}

\subsection{Center uniqueness}\label{sec:uniqueness}

A natural question is whether the semantic center is unique.  In
degenerate cases two coordinates may share the same propagation score,
but for ``generic'' weight assignments the center is unique.

\begin{definition}[Generic weights]\label{def:generic}
A weight function $\wt : E \to \mathbb{R}_{>0}$ is \emph{generic} if
the induced propagation scores $\{\PropScore(v)\}_{v \in V}$ are
pairwise distinct.
\end{definition}

\begin{theorem}[Center uniqueness]\label{thm:uniqueness}
If the weight function $\wt$ is generic, then the semantic center
$\SC(\Site)$ is the unique maximiser of $\PropScore$.
\end{theorem}
\begin{proof}
By definition, $\SC(\Site) = \operatorname{argmax}_v \PropScore(v)$.
If the propagation scores are pairwise distinct, the maximum is attained
at exactly one node.  Thus the tie-breaking rule in
Definition~\ref{def:semantic-center} is never invoked, and the center
is unique.
\end{proof}

\begin{remark}
The genericity condition holds for ``almost all'' weight assignments in
the measure-theoretic sense: the set of weight functions producing a
tie has Lebesgue measure zero in $\mathbb{R}_{>0}^{|E|}$.  In the
\JuGeo{} codebase, where weights derive from trust levels (a finite
but rich lattice) and morphism multiplicities, ties have not been
observed empirically.
\end{remark}

\subsection{Approximation bounds for DAGs}\label{sec:approx-dag}

When $\MG$ is a general DAG rather than a tree, verifying a coordinate
$v$ may reduce the effective weight of downstream nodes (their
remaining obligations shrink).  The weight function therefore becomes
\emph{dynamic}, and the optimal ordering is NP-hard to compute in
general (it reduces to weighted minimum completion time with
precedence constraints).  However, the greedy center-first heuristic
provides a bounded approximation.

\begin{definition}[Dynamic weight function]\label{def:dyn-weight}
Let $S_k \subseteq V$ denote the set of coordinates verified by step
$k$.  The \emph{residual weight} of $v \notin S_k$ is
\[
  \wt_k(v) = |\DepSet(v) \setminus S_k|,
\]
\ie{} only unverified successors contribute.
\end{definition}

\begin{theorem}[Greedy approximation ratio]\label{thm:approx}
Let $\anyord_{\mathrm{greedy}}$ be the ordering that at each step
verifies the coordinate with the largest residual weight.  For any DAG
with $n$ nodes,
\[
  \wtotal(\anyord_{\mathrm{greedy}})
  \leq (1 + \ln n) \cdot \wtotal(\sortord),
\]
where $\sortord$ is the optimal ordering.
\end{theorem}
\begin{proof}[Proof sketch]
The proof adapts the analysis of the greedy algorithm for weighted
set cover.  At each step the greedy choice picks the node $v^*$
maximising the ``cost reduction per unit'':
\[
  \Delta(v^*) = \max_{v \notin S_k}
  \frac{\wt_k(v)}{1}
  = \max_{v \notin S_k} \wt_k(v).
\]
Each unverified coordinate $u$ is ``covered'' when it enters $S_k$
(either directly or because all its predecessors in the DAG are
verified).  The harmonic-sum argument from the set-cover analysis
yields the $(1 + \ln n)$ factor: in the worst case, the
$j$-th most expensive coordinate contributes $\wt(v_j) / j$ to the
gap between greedy and optimal cost, and
$\sum_{j=1}^{n} 1/j \leq 1 + \ln n$.
\end{proof}

\begin{corollary}\label{cor:practical-gap}
For program dependency graphs with bounded out-degree $\Delta$, the
residual weights decrease by at least one per step, so the effective
approximation ratio is at most $1 + \ln \Delta$, which is typically
small ($\Delta \leq 20$ in the \JuGeo{} codebase, giving a ratio
$\leq 4.0$).
\end{corollary}

\subsection{Formal proof in Lean 4}

The theorem and its key lemma are formalised in
\LEAN{} in \texttt{proofs/lean/Paper50\_SemanticCenters.lean}.
The formalisation uses:
\begin{itemize}[noitemsep]
  \item \texttt{VerifNode} (a \texttt{structure} carrying an \texttt{id}
    and a \texttt{weight}),
  \item \texttt{totalCostFrom} (a recursive function computing the
    weighted sum starting from a given position),
  \item \texttt{swap\_reduces\_cost} (the pairwise exchange lemma),
  \item \texttt{center\_optimal\_two} (the 2-node base case), and
  \item \texttt{greedyCenterOptimal} (the full optimality statement for
    a sorted ordering).
\end{itemize}
All proofs are closed by \texttt{omega} or direct induction on list
structure; no \texttt{sorry} is used.

% =====================================================================
\section{Experiments}\label{sec:experiments}
% =====================================================================

\subsection{Setup}

We evaluate center-first verification on the \JuGeo{} codebase
(approximately $350$ verification coordinates).  Three orderings are
compared:
\begin{description}[noitemsep]
  \item[Random] A uniformly random permutation of coordinates.
  \item[Alphabetical] Coordinates sorted lexicographically.
  \item[Center-first (\SCD{})] Coordinates ordered by the output of
        Algorithm~\ref{alg:scd}.
\end{description}
Each run verifies all $350$ coordinates using the Z3 SMT solver
(Paper~03) with a per-coordinate timeout of $5\,\mathrm{s}$.
We measure total wall-clock time and the fraction of coordinates
reaching trust $\geq \Tcopilot$ within each checkpoint.

\subsection{Results}

\begin{table}[h]
\centering
\caption{Verification performance under three orderings.
         Mean $\pm$ std.\ dev.\ over 10 runs.}
\label{tab:results}
\begin{tabularx}{\textwidth}{Xrrr}
\toprule
\textbf{Ordering} &
\textbf{Total time (s)} &
\textbf{Steps to 50\% trust} &
\textbf{Steps to 87\% trust} \\
\midrule
Random              & \ppFiftyRandomTime & \ppFiftyRandomStepsHalf & \ppFiftyRandomStepsFull \\
Alphabetical        & \ppFiftyAlphaTime & \ppFiftyAlphaStepsHalf & \ppFiftyAlphaStepsFull \\
Center-first (\SCD{}) & \ppFiftyCenterTime & \ppFiftyCenterStepsHalf & \ppFiftyCenterStepsFull \\
\bottomrule
\end{tabularx}
\end{table}

Center-first achieves a significant reduction in total
verification time (Center Optimality Theorem) and reaches trust
coverage faster than alternatives (universal benchmark mean \expTimeMean{}).

\subsection{Trust propagation curve}

Figure~\ref{fig:trust-curve} (described textually) shows the trust
coverage $\VProp(k)$ (fraction of coordinates with $\trust \geq
\Tcopilot$) as a function of verification steps $k$:

\begin{itemize}[noitemsep]
  \item \textbf{Center-first}: $\VProp$ grows steeply in the first
    $50$ steps (trust propagates outward from \texttt{Judgment}),
    then levels off as leaf coordinates are addressed.
  \item \textbf{Alphabetical}: $\VProp$ grows roughly linearly.
  \item \textbf{Random}: $\VProp$ follows a sub-linear S-curve due to
    occasional lucky early hits.
\end{itemize}

The area under the trust-coverage curve (AUTC) is highest for
center-first, confirming the theoretical optimality,
reflecting that center-first maintains high trust coverage throughout
the verification run.

\subsection{Sensitivity to $\alpha$}

We vary the mixing parameter $\alpha \in \{0, 0.25, 0.5, 0.75, 1\}$
controlling the weight between betweenness centrality and propagation
rank.  Total verification time is minimised near
$\alpha = 0.5$ and increases at the extremes ($\alpha = 0$ or $\alpha = 1$),
suggesting that the composite score is more effective than either
measure alone.  The default $\alpha = 0.5$ is robust.

\subsection{Comparison with semantic caching}

Paper~38 showed that semantic caching reduces
latency.  Combining caching with center-first ordering
yields further improvement: the cache serves calls at
coordinates that were already informed by the center, requiring no
solver invocation (site success rate \subSiteSuccessRate, mean \subSiteMeanTime).

% =====================================================================
\section{Conclusion}\label{sec:conclusion}
% =====================================================================

We have introduced \emph{semantic centers}---coordinates that maximise
propagation score in the morphism graph of a semantic site---and proved
that verifying them first minimises total weighted verification cost for
tree-structured sites.  The \texttt{SemanticCenterDetectionAlgorithm}
detects the center in linear time on sparse graphs, and the
\texttt{ThesisClaimTracker} propagates trust outward through the
morphism graph as verification proceeds.  Empirically, center-first
ordering delivers a significant reduction in verification time on the
\JuGeo{} codebase (universal benchmark: \expTotalPrograms{} programs,
\expAccuracy{} accuracy, mean \expTimeMean).

\paragraph{Future work.}
\begin{itemize}[noitemsep]
  \item \emph{Dynamic weights.}  Extend the optimality theorem to DAGs
    where weights change as information propagates
    (requires dynamic programming or approximation).
  \item \emph{Multi-center strategies.}  Identify a small
    \emph{core set} of coordinates whose joint propagation covers
    the site.
  \item \emph{Integration with semantic caching.}  Use the center
    ordering to pre-warm the semantic cache before a verification run.
  \item \emph{Trust-sensitive weighting.}  Use the trust level of
    candidate judgments as a multiplier on $\wt(v)$, so that
    well-evidenced nodes are preferred over weakly-evidenced ones of
    equal degree.
\end{itemize}

\paragraph{Acknowledgements.}
The semantic center concept was inspired by the betweenness-centrality
literature in network science~\cite{Freeman1977} and the weighted-job
scheduling results of~\cite{Conway1967}.

\begin{thebibliography}{99}

\bibitem{jg01}
H.~Young.
\newblock {Judgment Geometry: A Semantic Site Framework for Program
  Verification}.
\newblock Paper~01 of the Judgment Geometry series, 2024.

\bibitem{jg04}
H.~Young.
\newblock {The Trust Algebra of Semantic Judgments}.
\newblock Paper~04 of the Judgment Geometry series, 2024.

\bibitem{jg06}
H.~Young.
\newblock {Semantic Moves: Navigation in Judgment Space}.
\newblock Paper~06 of the Judgment Geometry series, 2024.

\bibitem{jg38}
H.~Young.
\newblock {Semantic Caching and Incremental Verification in the Judgment
  Geometry Framework}.
\newblock Paper~38 of the Judgment Geometry series, 2025.

\bibitem{Brandes2001}
U.~Brandes.
\newblock A faster algorithm for betweenness centrality.
\newblock \emph{Journal of Mathematical Sociology}, 25(2):163--177, 2001.

\bibitem{Conway1967}
R.~W.~Conway, W.~L.~Maxwell, and L.~W.~Miller.
\newblock \emph{Theory of Scheduling}.
\newblock Addison-Wesley, 1967.

\bibitem{Freeman1977}
L.~C.~Freeman.
\newblock A set of measures of centrality based on betweenness.
\newblock \emph{Sociometry}, 40(1):35--41, 1977.

\bibitem{Page1999}
L.~Page, S.~Brin, R.~Motwani, and T.~Winograd.
\newblock The {PageRank} citation ranking: Bringing order to the web.
\newblock Technical Report 1999-66, Stanford InfoLab, 1999.

\bibitem{Graham1969}
R.~L.~Graham.
\newblock Bounds on multiprocessing timing anomalies.
\newblock \emph{SIAM Journal on Applied Mathematics}, 17(2):416--429, 1969.

\end{thebibliography}

\end{document}
